% ##############################################################################
\section{State of the art: industry landscape}
\label{sec:state_of_the_art}

Section~\ref{sec:related_work} positions \Noesis{} against academic work on
routing, fusion, distillation, and auction theory. This section instead
surveys currently operating startups and projects that address, piece by
piece, the same problem \Noesis{} targets: buying and selling machine
intelligence, trading raw compute, or verifying that a delivered service met
its promise. The goal is to locate the specific combination of properties,
auction-based price discovery for a capability/latency/reliability bundle,
metered fulfillment, and reputation feedback into pricing, that no surveyed
system currently occupies. The industry moves quickly, so the picture below
reflects the state of the market at the time of writing and should be read
as illustrative of the categories rather than exhaustive within any one of
them.

% ==============================================================================
\subsection{Posted-price multi-provider gateways}
\label{sec:sota_gateways}

OpenRouter, introduced in Section~\ref{sec:protocol_overview}, is one of
several gateways that unify access to hundreds of models behind a single
API and bill at or near the underlying provider's price. Portkey,
LiteLLM, Requesty, TrueFoundry's AI Gateway, Cloudflare AI Gateway, and the
Vercel AI Gateway occupy the same niche, adding a fixed platform fee, a
flat markup, or no markup at all on top of pass-through provider pricing,
together with observability, caching, and automatic failover across
providers. A single-vendor gateway such as AWS Bedrock offers the same
unified-API convenience without the multi-seller pooling. OpenRouter
additionally ships a mixture-of-agents feature, Fusion, that fans a prompt
out to several models and synthesizes one answer, which is the closest
commercial analog to the answer-fusion strategy of
Section~\ref{sec:fusion} that this survey identified.
Table~\ref{tab:sota_gateways_companies} lists the systems in this cluster.

\begin{samepage}
  \begin{table}[H]
    \centering
    \footnotesize
    \begin{tabular}{l|l|p{5.8cm}|}
                                & \emph{URL} & \emph{Distinguishing feature} \\
      \hline
      OpenRouter                & \url{https://openrouter.ai} & Multi-vendor pooling; ships Fusion mixture-of-agents synthesis \\
      \hline
      Portkey                   & \url{https://portkey.ai} & Pass-through pricing + platform fee/markup$^\dagger$ \\
      \hline
      LiteLLM                   & \url{https://litellm.ai} & Pass-through pricing + platform fee/markup$^\dagger$ \\
      \hline
      Requesty                  & \url{https://requesty.ai} & Pass-through pricing + platform fee/markup$^\dagger$ \\
      \hline
      TrueFoundry AI Gateway     & \url{https://truefoundry.com} & Pass-through pricing + platform fee/markup$^\dagger$ \\
      \hline
      Cloudflare AI Gateway      & \url{https://developers.cloudflare.com/ai-gateway} & Pass-through pricing + platform fee/markup$^\dagger$ \\
      \hline
      Vercel AI Gateway          & \url{https://vercel.com/ai-gateway} & Pass-through pricing + platform fee/markup$^\dagger$ \\
      \hline
      AWS Bedrock                & \url{https://aws.amazon.com/bedrock} & Single-vendor; unified API without multi-seller pooling \\
      \hline
    \end{tabular}
    \caption{Posted-price multi-provider gateways surveyed in
      Section~\ref{sec:sota_gateways}. $^\dagger$Also bundles
      observability, caching, and automatic failover across providers.}
    \label{tab:sota_gateways_companies}
  \end{table}
\end{samepage}

None of these gateways clears an auction: every price is posted or passed
through unilaterally, so there is no price discovery in the sense of
Section~\ref{sec:comparison_table}. None meters whether a request's
capability, latency, or reliability matched a stated contract beyond a
standard infrastructure uptime SLA, and none turns a delivered shortfall
into a reputation signal that changes a provider's standing in future
transactions.

% ==============================================================================
\subsection{Learned model routers}
\label{sec:sota_routing}

A second cluster of startups routes each individual request to the model
expected to give the best cost/quality trade-off, which is the commercial
counterpart to the difficulty-aware routing of Section~\ref{sec:routing}.
Martian routes on an interpretability-derived model of what each request
needs. Not Diamond and Unify learn or benchmark a routing policy across
models on standard evaluation suites. Arch, an open-source, Envoy-based
prompt gateway, routes from a natural-language policy description rather
than hardcoded rules. Table~\ref{tab:sota_routing_companies} lists the
systems in this cluster.

\begin{samepage}
  \begin{table}[H]
    \centering
    \footnotesize
    \begin{tabular}{l|l|p{5.8cm}|}
                    & \emph{URL} & \emph{Distinguishing feature} \\
      \hline
      Martian        & \url{https://withmartian.com} & Routes on an interpretability-derived model of request needs \\
      \hline
      Not Diamond    & \url{https://notdiamond.ai} & Learns/benchmarks a routing policy on standard evaluation suites \\
      \hline
      Unify          & \url{https://unify.ai} & Learns/benchmarks a routing policy on standard evaluation suites \\
      \hline
      Arch           & \url{https://github.com/katanemo/archgw} & Open-source, Envoy-based; routes from a natural-language policy description \\
      \hline
    \end{tabular}
    \caption{Learned model routers surveyed in Section~\ref{sec:sota_routing}.}
    \label{tab:sota_routing_companies}
  \end{table}
\end{samepage}

These routers make the same difficulty-aware bet as \NoesisServer, that most
requests do not need the most capable model, but the routing decision is a
private policy internal to each vendor rather than the outcome of a market
that prices the tiers being routed between, and none of them exposes a
reputation signal that a request's actual capability, latency, or
reliability fell short of what was routed for.

% ==============================================================================
\subsection{Reverse-auction and spot compute marketplaces}
\label{sec:sota_compute_markets}

A third cluster trades raw compute, rather than model access, through
genuine bidding mechanisms. Akash Network runs a reverse auction in which
a tenant posts a workload specification and a budget and providers bid to
fill it. Vast.ai runs a continuous bid/ask market for GPU-hours in which
an underbid spot instance can be reclaimed by a higher bidder. RunPod,
io.net, Golem Network, CoreWeave, and Prime Intellect offer spot, reserved,
or brokered tiers over pooled GPU capacity, with pricing that ranges from
genuinely biddable to posted with a spot discount. Table~\ref{tab:sota_compute_markets_companies}
lists the systems in this cluster.

\begin{samepage}
  \begin{table}[H]
    \centering
    \footnotesize
    \begin{tabular}{l|l|c|p{4.3cm}|}
                       & \emph{URL} & \emph{Bid-based} & \emph{Distinguishing feature} \\
      \hline
      Akash Network     & \url{https://akash.network} & \ngood & Reverse auction: tenant posts workload + budget, providers bid \\
      \hline
      Vast.ai           & \url{https://vast.ai} & \ngood & Continuous bid/ask GPU-hour market; underbid spot instances reclaimable \\
      \hline
      RunPod            & \url{https://runpod.io} & \npartial & Spot/reserved/brokered GPU tiers$^\dagger$ \\
      \hline
      io.net            & \url{https://io.net} & \npartial & Spot/reserved/brokered GPU tiers$^\dagger$ \\
      \hline
      Golem Network     & \url{https://golem.network} & \npartial & Spot/reserved/brokered GPU tiers$^\dagger$ \\
      \hline
      CoreWeave         & \url{https://coreweave.com} & \npartial & Spot/reserved/brokered GPU tiers$^\dagger$ \\
      \hline
      Prime Intellect    & \url{https://primeintellect.ai} & \npartial & Spot/reserved/brokered GPU tiers$^\dagger$ \\
      \hline
    \end{tabular}
    \caption{Reverse-auction and spot compute marketplaces surveyed in
      Section~\ref{sec:sota_compute_markets}. $^\dagger$Pricing ranges from
      genuinely biddable to posted with a spot discount.}
    \label{tab:sota_compute_markets_companies}
  \end{table}
\end{samepage}

Akash and Vast.ai in particular demonstrate that an open bid/ask market for
compute-time is viable at production scale, which is direct evidence for
the auction half of \NoesisMarket's design. The traded unit, however, is a
GPU-hour or a workload's raw resource footprint, not the (capability,
latency, reliability) bundle of Section~\ref{sec:intelligence_as_commodity}:
none of these markets prices, or even observes, how good the resulting
model answer is, so the bundled-quality and difficulty-routing rows of
Table~\ref{tab:comparison} remain unaddressed.

% ==============================================================================
\subsection{Token-incentivized inference networks}
\label{sec:sota_decentralized}

A fourth cluster builds reputation and verification directly into a
decentralized compute network, which is the closest analog to the
fulfillment-monitoring and reputation loop that connects \NoesisServer{}
and \NoesisMarket{} (Sections~\ref{sec:fulfillment_monitoring} and
\ref{sec:reputation}).
Bittensor organizes miners and validators into subnets; in its
inference-oriented subnets, validators send prompts to many miners, score
the returned completions for quality and latency, and an on-chain
consensus mechanism aggregates the stake-weighted scores into recurring
reward emissions, with a serverless-inference subnet layered on top for
pay-per-query access. Gensyn verifies distributed training work through a
probabilistic proof-of-learning protocol rather than scoring inference
output. Ritual runs a network of inference oracles that attach a
cryptographic proof of correct execution to each computation. Nosana and
the Fetch.ai / ASI Alliance pool idle compute and agent-to-agent
transactions for AI workloads more broadly.
Table~\ref{tab:sota_decentralized_companies} lists the systems in this
cluster.

\begin{samepage}
  \begin{table}[H]
    \centering
    \footnotesize
    \begin{tabular}{l|l|p{6.2cm}|}
                             & \emph{URL} & \emph{Distinguishing feature} \\
      \hline
      Bittensor               & \url{https://bittensor.com} & Validator scoring of inference quality/latency; stake-weighted reward emissions \\
      \hline
      Gensyn                  & \url{https://gensyn.ai} & Probabilistic proof-of-learning verifies distributed training \\
      \hline
      Ritual                  & \url{https://ritual.net} & Inference oracles attach a cryptographic proof of correct execution \\
      \hline
      Nosana                  & \url{https://nosana.com} & Pools idle compute and agent-to-agent transactions for AI workloads \\
      \hline
      Fetch.ai / ASI Alliance  & \url{https://fetch.ai} & Pools idle compute and agent-to-agent transactions for AI workloads \\
      \hline
    \end{tabular}
    \caption{Token-incentivized inference networks surveyed in
      Section~\ref{sec:sota_decentralized}.}
    \label{tab:sota_decentralized_companies}
  \end{table}
\end{samepage}

Bittensor's validator scoring is the only mechanism surveyed here that
continuously measures output quality and lets that measurement change a
provider's economic standing, which is precisely the loop
Section~\ref{sec:reputation} formalizes for \NoesisMarket. It differs from
\Noesis{} in two respects: the reward is an emission schedule internal to
the network's consensus rather than a price cleared against an explicit
buyer bid, and the score is a continuous internal ranking rather than a
pass/fail check against a specific bought (volume, tier, latency,
reliability) contract. Gensyn's and Ritual's proofs verify that a
computation was executed correctly, not that its output met a stated
capability or latency bound, so they address a different, complementary
notion of trust.

% ==============================================================================
\subsection{Compute exchanges and derivatives}
\label{sec:sota_exchanges}

A fifth, more recent cluster builds explicit auction or index
infrastructure on top of the raw-compute markets of
Section~\ref{sec:sota_compute_markets}. SF Compute operates an
order-book-style market for GPU-time blocks. Compute Exchange runs spot
auctions over standardized compute contracts explicitly modeled on a
commodity exchange. Andromeda brokers capacity across many providers on
standardized, benchmarked contracts to avoid multi-year lock-in. Ornn's
Compute Price Index publishes a transaction-based benchmark price for GPU
rental, feeding into early cash-settled compute futures.
Table~\ref{tab:sota_exchanges_companies} lists the systems in this
cluster.

\begin{samepage}
  \begin{table}[H]
    \centering
    \footnotesize
    \begin{tabular}{l|l|p{5.8cm}|}
                        & \emph{URL} & \emph{Distinguishing feature} \\
      \hline
      SF Compute         & \url{https://sfcompute.com} & Order-book-style market for GPU-time blocks \\
      \hline
      Compute Exchange   & \url{https://compute.exchange} & Spot auctions over standardized compute contracts \\
      \hline
      Andromeda          & \url{https://andromeda.ai} & Brokers capacity across providers on standardized, benchmarked contracts \\
      \hline
      Ornn               & \url{https://ornn.ai} & Compute Price Index: transaction-based benchmark price feeding cash-settled futures \\
      \hline
    \end{tabular}
    \caption{Compute exchanges and derivatives surveyed in
      Section~\ref{sec:sota_exchanges}.}
    \label{tab:sota_exchanges_companies}
  \end{table}
\end{samepage}

This cluster is the strongest evidence that the pairing of an open auction
with a standardized contract, the structural bet underlying
\NoesisMarket, is a fundable design for AI infrastructure: it did not exist
a few years ago and now spans spot auctions, brokered standardized
contracts, and price indices. Every instance found, however, prices raw
compute capacity rather than a verified LLM answer, so none of them closes
the fulfillment-to-reputation loop that Section~\ref{sec:noesis_server}
adds on top of the market.

% ==============================================================================
\subsection{Serverless inference providers as prospective sellers}
\label{sec:sota_sellers}

A final cluster, serverless inference providers such as Together AI,
Fireworks AI, Groq, Baseten, and Fal.ai (Table~\ref{tab:sota_sellers_companies}),
sell inference at posted per-token or per-second rates at a given
capability tier. These are not competitors to \Noesis{} so much as its
prospective sellers: each already meters its own delivered latency and
offers a standard infrastructure SLA, which is a useful starting point for
the request-level fulfillment statistics of
Section~\ref{sec:fulfillment_monitoring}, but none of them exposes that
metering as a contract a buyer can shop across sellers for, or as a
reputation signal that follows the seller into a future transaction.

\begin{samepage}
  \begin{table}[H]
    \centering
    \footnotesize
    \begin{tabular}{l|l|}
                & \emph{URL} \\
      \hline
      Together AI  & \url{https://together.ai} \\
      \hline
      Fireworks AI & \url{https://fireworks.ai} \\
      \hline
      Groq         & \url{https://groq.com} \\
      \hline
      Baseten      & \url{https://baseten.co} \\
      \hline
      Fal.ai       & \url{https://fal.ai} \\
      \hline
    \end{tabular}
    \caption{Serverless inference providers surveyed in
      Section~\ref{sec:sota_sellers}, each selling inference at posted
      per-token or per-second rates at a given capability tier.}
    \label{tab:sota_sellers_companies}
  \end{table}
\end{samepage}

% ==============================================================================
\subsection{Comparison and market gap}
\label{sec:sota_gap}

Table~\ref{tab:sota_comparison} extends the comparison of
Table~\ref{tab:comparison} to the six clusters surveyed above, each
represented by its most characteristic mechanism.

Figure~\ref{fig:sota_landscape} plots the six clusters on these two
axes.

% rendered_images:begin
% ```tikz
% \definecolor{ptcolor}{RGB}{74,110,168}
% \definecolor{noesiscolor}{RGB}{178,58,72}
% \draw[->, >=stealth, thick] (0,0) -- (9.6,0) node[right, font=\footnotesize, align=left] {auction-based\\price discovery};
% \draw[->, >=stealth, thick] (0,0) -- (0,7.4) node[above, font=\footnotesize, align=center] {bundled, metered\\quality};
% \draw[dashed, gray] (4.6,-0.3) -- (4.6,7.4);
% \draw[dashed, gray] (-0.3,3.6) -- (9.6,3.6);
% \node[font=\footnotesize] at (2.3,-0.5) {no};
% \node[font=\footnotesize] at (7.0,-0.5) {yes};
% \node[font=\footnotesize] at (-0.6,1.8) {no};
% \node[font=\footnotesize] at (-0.6,5.5) {yes};
% \filldraw[ptcolor] (0.9,0.9) circle (2.4pt);
% \node[anchor=west, align=left, font=\footnotesize] at (1.3,0.9) {learned routers\\(Martian, Not Diamond,\\Unify, Arch)};
% \filldraw[ptcolor] (0.9,2.4) circle (2.4pt);
% \node[anchor=west, align=left, font=\footnotesize] at (1.3,2.4) {posted-price gateways\\(OpenRouter, Portkey,\\LiteLLM, ...)};
% \filldraw[ptcolor] (0.9,3.3) circle (2.4pt);
% \node[anchor=west, align=left, font=\footnotesize] at (1.3,3.3) {serverless inference providers\\(Together, Fireworks, Groq, ...)};
% \filldraw[ptcolor] (0.9,4.9) circle (2.4pt);
% \node[anchor=west, align=left, font=\footnotesize] at (1.3,4.9) {token-incentivized\\inference networks\\(Bittensor, Gensyn, Ritual)};
% \filldraw[ptcolor] (5.2,0.9) circle (2.4pt);
% \node[anchor=west, align=left, font=\footnotesize] at (5.6,0.9) {reverse-auction \& spot compute\\(Akash, Vast.ai, RunPod, ...)};
% \filldraw[ptcolor] (5.2,2.1) circle (2.4pt);
% \node[anchor=west, align=left, font=\footnotesize] at (5.6,2.1) {compute exchanges \& derivatives\\(SF Compute, Compute Exchange, ...)};
% \filldraw[noesiscolor] (8.8,6.6) circle (3pt);
% \node[anchor=west, align=left, font=\footnotesize\bfseries, text=noesiscolor] at (8.3,7.05) {Noesis};
% ```
% label=fig:sota_landscape
% caption=The six industry clusters of Table~\ref{tab:sota_comparison}
%   positioned by auction-based price discovery and bundled, metered
%   quality. Gateways, learned routers, and serverless providers set
%   prices unilaterally; compute marketplaces and exchanges bid on raw
%   compute but neither prices nor observes answer quality;
%   token-incentivized networks verify quality without a market clearing
%   price against a buyer's bid. No surveyed cluster combines both
%   properties for machine intelligence; \Noesis{} occupies that
%   combination.
% rendered_images:end
% render_images:begin
\begin{figure}[H]
  \includegraphics[width=\linewidth]{07_state_of_the_art.tex.figs/07_state_of_the_art.1.png}
  \caption{The six industry clusters of Table~\ref{tab:sota_comparison} positioned by auction-based price discovery and bundled, metered quality. Gateways, learned routers, and serverless providers set prices unilaterally; compute marketplaces and exchanges bid on raw compute but neither prices nor observes answer quality; token-incentivized networks verify quality without a market clearing price against a buyer's bid. No surveyed cluster combines both properties for machine intelligence; \Noesis{} occupies that combination.}
  \label{fig:sota_landscape}
\end{figure}
% render_images:end

\begin{samepage}
  \begin{table}[H]
    \centering
    \footnotesize
    \begin{tabular}{l|c|c|c|c|c|}
                                     & \emph{Auction} & \emph{Bundled} & \emph{Metered} & \emph{Reputation} & \emph{Liquidity} \\
                                     & \emph{price}   & \emph{quality} & \emph{fulfillment} & \emph{$\to$ pricing} & \emph{pooling} \\
      \hline
      Posted-price gateways (\S\ref{sec:sota_gateways})            & \nbad & \npartial & \nbad     & \nbad     & \ngood \\
      \hline
      Learned model routers (\S\ref{sec:sota_routing})             & \nbad & \npartial & \nbad     & \nbad     & \npartial \\
      \hline
      Reverse-auction compute markets (\S\ref{sec:sota_compute_markets}) & \ngood & \nbad     & \nbad     & \npartial & \ngood \\
      \hline
      Token-incentivized inference networks (\S\ref{sec:sota_decentralized}) & \nbad & \nbad     & \npartial & \ngood    & \npartial \\
      \hline
      Compute exchanges \& derivatives (\S\ref{sec:sota_exchanges}) & \ngood & \nbad     & \npartial & \npartial & \ngood \\
      \hline
      Serverless inference sellers (\S\ref{sec:sota_sellers})       & \nbad & \npartial & \nbad     & \nbad     & \nbad  \\
      \hline
      \Noesis                                                       & \ngood & \ngood    & \ngood    & \ngood    & \ngood \\
      \hline
    \end{tabular}
    \caption{Comparison of \Noesis{} against the six industry clusters
      surveyed in Section~\ref{sec:state_of_the_art}. ``Bundled quality''
      denotes a purchasable (capability, latency, reliability) contract
      rather than a raw compute quantity or a single posted price;
      ``metered fulfillment'' denotes verification of delivered quality
      against that contract, rather than a generic infrastructure SLA.}
    \label{tab:sota_comparison}
  \end{table}
\end{samepage}

No cluster combines auction-based price discovery with a bundled,
verifiable quality contract and a reputation loop that feeds back into
pricing. What is striking is that each individual ingredient is already
validated in production by at least one system: Akash and Vast.ai show
that an open auction for compute is viable; Compute Exchange and Andromeda
show that standardizing the traded contract attracts liquidity; Bittensor
shows that continuously scoring delivered quality and tying it to economic
standing changes provider incentives; OpenRouter's Fusion feature shows
that answer fusion is a sellable product feature. \Noesis's contribution is
not any one of these mechanisms in isolation, each already has a
real-world precedent, but their combination into a single loop, applied to
a verifiable (capability, latency, reliability) bundle rather than to raw
compute, which to our knowledge no surveyed system currently offers.
